%% 致谢

\begin{acknowledgements}
在本课题研究与论文撰写过程中，指导教师在研究方向确定、任务体系设计、实验分析和论文结构梳理等方面给予了持续指导。老师对论文中心思想、实验结论深度和材料呈现方式提出的建议，使本文能够从单纯的结果整理进一步转向围绕无人机双认知能力展开的系统分析。

同时，感谢课题推进过程中为实验环境、模型调用、数据整理和材料检查提供帮助的同学与朋友。UAV-DualCog 的构建涉及场景扫描、地标复核、视频生成、图像任务生成、多模型实验和展示材料整理等多个环节，相关讨论和反馈帮助我及时发现并改进了许多工程细节问题。

最后，感谢家人在本科阶段学习和毕业设计期间给予的理解与支持。毕业设计不仅是一次技术实践，也是一次对研究问题、工程实现和学术表达能力的综合训练。谨向所有给予帮助和鼓励的人表示诚挚感谢。
\end{acknowledgements}
